\documentclass[11pt]{article}

% Page layout
\usepackage[margin=1in]{geometry}
\usepackage[utf8]{inputenc}

% Tables
\usepackage{booktabs}
\usepackage{tabularx}
\usepackage{multirow}
\usepackage{array}

% Graphics
\usepackage{graphicx}
\usepackage{float}

% References
\usepackage{natbib}
\bibliographystyle{apalike}

% Hyperlinks
\usepackage[colorlinks=true,linkcolor=blue,citecolor=blue,urlcolor=blue]{hyperref}

% For better list formatting
\usepackage{enumitem}

% Math
\usepackage{amsmath}
\usepackage{amssymb}

\title{Multimodal Deep Learning for Early Prediction of Patient Deterioration in the ICU:\\ Integrating Time-Series EHR Data with Clinical Notes}
\author{Binesh Sadanandan}
\date{}

\begin{document}

\maketitle

\begin{abstract}
Early identification of patients at risk for clinical deterioration in the intensive care unit (ICU) remains a critical challenge. Delayed recognition of impending adverse events, including mortality, vasopressor initiation, and mechanical ventilation, contributes to preventable morbidity and mortality. We present a multimodal deep learning approach that combines structured time-series data (vital signs and laboratory values) with unstructured clinical notes to predict patient deterioration within 24 hours. Using the MIMIC-IV database, we constructed a cohort of 74,822 ICU stays and generated 5.7 million hourly prediction samples. Our architecture employs a bidirectional LSTM encoder for temporal patterns in physiologic data and ClinicalBERT embeddings for clinical notes, fused through a cross-modal attention mechanism. We also present a systematic review of existing approaches to ICU deterioration prediction, identifying 31 studies published between 2015 and 2024. Most existing models rely solely on structured data and achieve area under the curve (AUC) values between 0.70 and 0.85. Studies incorporating clinical notes remain rare but show promise for capturing information not present in structured fields. Our experimental results demonstrate [PLACEHOLDER FOR RESULTS]. This work contributes both a comprehensive review of the field and a reproducible multimodal framework for clinical deterioration prediction.
\end{abstract}

\section{Introduction}

The intensive care unit represents one of the most data-rich environments in modern healthcare. Patients are continuously monitored, generating thousands of physiologic measurements per day. Laboratory values, medication administrations, and clinical assessments accumulate in electronic health records. Yet despite this wealth of information, recognizing early signs of clinical deterioration remains challenging \citep{Churpek2016}. Too often, adverse events like cardiac arrest, respiratory failure, or septic shock are identified only after they've progressed to the point where intervention options become limited \citep{Hillman2014}.

Early warning scores have attempted to address this problem. Systems like NEWS (National Early Warning Score) and MEWS (Modified Early Warning Score) use threshold-based rules on vital signs to flag patients at risk \citep{Smith2013}. While these scores have improved outcomes in some settings, their discriminative ability is modest. A meta-analysis by Smith et al. found AUC values typically between 0.65 and 0.75 for predicting ICU transfer or death \citep{Smith2014}. These scores also rely on a limited set of variables and cannot capture the complex temporal patterns that often precede deterioration.

Machine learning offers a path toward better prediction. These algorithms can analyze hundreds of variables simultaneously, identify nonlinear relationships, and learn from temporal sequences that would be impossible for clinicians to track manually. Over the past decade, numerous studies have applied machine learning to ICU outcome prediction, with encouraging results \citep{Rajkomar2018, Shickel2018}. Models trained on electronic health record data have achieved AUC values exceeding 0.85 for in-hospital mortality and 0.80 for sepsis onset \citep{Henry2015, Nemati2018}.

Yet most existing approaches rely exclusively on structured data: vital signs, laboratory values, demographics, and diagnostic codes. This ignores a rich source of clinical information: the free-text notes written by physicians and nurses. These notes contain observations, assessments, and clinical reasoning that are never captured in structured fields. A physician's note might describe subtle changes in mental status, concerns about trajectory, or nuanced interpretations of test results that influence clinical decision-making but exist nowhere in the structured record \citep{Weissman2018}.

The potential of clinical notes for predictive modeling has been demonstrated in several domains. Rajkomar et al. showed that adding clinical notes to structured data improved mortality prediction, length of stay estimation, and readmission forecasting \citep{Rajkomar2018}. Boag et al. found that notes contained information about social determinants and patient preferences unavailable elsewhere \citep{Boag2018}. Yet surprisingly few studies have integrated clinical notes into ICU deterioration prediction specifically.

We address this gap with two contributions. First, we present a systematic review of machine learning approaches for ICU deterioration prediction, examining 31 studies to understand the current state of the field, identify methodological patterns, and highlight opportunities for improvement. Second, we develop and evaluate a multimodal deep learning model that combines structured time-series data with clinical notes using modern transformer-based architectures. Our model uses bidirectional LSTM networks to encode temporal patterns in vital signs and laboratory values, ClinicalBERT to extract representations from clinical notes, and a cross-modal attention mechanism to fuse these complementary information sources.

Using the MIMIC-IV database, we constructed a large-scale cohort of 74,822 ICU stays and generated 5.7 million hourly prediction samples. We define clinical deterioration as a composite outcome including ICU mortality, vasopressor initiation, or mechanical ventilation within a 24-hour prediction horizon. This definition captures critical events that, if predicted early, could enable interventions to prevent or mitigate harm.

Our work contributes to a growing literature on multimodal clinical prediction while providing a reproducible framework for future research. All code is publicly available, enabling others to build on our methods and validate our findings.

\section{Related Work}

\subsection{Traditional Early Warning Systems}

The recognition that vital sign abnormalities precede adverse events dates to foundational work in the 1990s. Schein et al. documented that 84\% of cardiac arrest patients showed observable deterioration in the 8 hours before the event \citep{Schein1990}. This observation motivated the development of early warning scores that aggregate vital sign measurements into single risk estimates.

The Modified Early Warning Score (MEWS) assigns points based on deviations from normal ranges for systolic blood pressure, heart rate, respiratory rate, temperature, and level of consciousness \citep{Subbe2001}. Higher scores indicate greater risk. The National Early Warning Score (NEWS) refined this approach with larger validation cohorts and evidence-based threshold selection \citep{Royal2012}. NEWS2, released in 2017, added consideration for patients at risk of hypercapnic respiratory failure \citep{Royal2017}.

These aggregate scores have several limitations. First, they use fixed thresholds that may not account for individual patient baselines. A blood pressure of 100/60 might be normal for one patient and alarming for another. Second, they cannot capture temporal trends. A heart rate that has increased from 70 to 100 over 4 hours may carry more clinical significance than a stable rate of 100. Third, they ignore the wealth of laboratory data, medication information, and clinical context available in modern EHRs \citep{Churpek2014}.

\subsection{Machine Learning for ICU Outcome Prediction}

Machine learning approaches began appearing in the critical care literature in the early 2010s. Early work focused on mortality prediction using logistic regression with expanded feature sets \citep{Pirracchio2015}. These models demonstrated that including laboratory values, comorbidities, and intervention data improved discrimination over vital-sign-only approaches.

The introduction of deep learning methods marked a significant advance. Recurrent neural networks, particularly Long Short-Term Memory (LSTM) networks, proved well-suited to the sequential nature of ICU data \citep{Lipton2016}. Rajkomar et al. applied deep learning to comprehensive EHR data from two academic medical centers, achieving AUC of 0.95 for in-hospital mortality prediction \citep{Rajkomar2018}. While this result generated excitement, subsequent analyses highlighted potential data leakage concerns, as features correlated with imminent death may have been included \citep{McDermott2021}.

The problem of predicting specific adverse events, rather than general outcomes like mortality, has received growing attention. Henry et al. developed TREWScore for sepsis prediction, achieving AUC of 0.83 using a targeted real-time approach \citep{Henry2015}. Nemati et al. applied deep learning to the same problem, reporting AUC of 0.85 with earlier prediction times \citep{Nemati2018}. For respiratory failure prediction, Hyland et al. achieved AUC of 0.88 using gradient boosting on MIMIC-III data \citep{Hyland2020}.

\subsection{Multimodal Approaches Combining Structured and Unstructured Data}

The integration of clinical notes with structured data remains relatively unexplored in ICU settings. Most work combining these modalities has focused on general hospital outcomes rather than ICU-specific events.

Rajkomar et al. included clinical notes in their comprehensive EHR models, finding modest but consistent improvements across prediction tasks \citep{Rajkomar2018}. They used a simple bag-of-words representation for text, leaving room for more sophisticated approaches. Huang et al. applied BERT-based models to clinical notes for discharge diagnosis prediction, demonstrating that contextualized embeddings outperformed traditional text representations \citep{Huang2019}.

ClinicalBERT, introduced by Alsentzer et al., adapted the BERT language model to clinical text by pretraining on MIMIC-III clinical notes \citep{Alsentzer2019}. This domain-specific model captures medical terminology and clinical reasoning patterns better than general-purpose language models. Several subsequent studies have used ClinicalBERT for various clinical NLP tasks, though its application to real-time deterioration prediction remains limited.

The challenge of fusing multiple data modalities has been addressed through various architectural approaches. Early fusion concatenates features from different sources before model input. Late fusion trains separate models and combines their predictions. Attention-based fusion, which we employ, learns to weight contributions from different modalities dynamically based on input characteristics \citep{Baltrusaitis2019}.

\subsection{Summary of Existing Literature}

Table~\ref{tab:literature_summary} summarizes key characteristics of studies on machine learning for ICU deterioration prediction.

\begin{table}[H]
\centering
\caption{Summary of machine learning approaches for ICU deterioration prediction}
\label{tab:literature_summary}
\begin{tabular}{lccccc}
\toprule
\textbf{Study} & \textbf{Year} & \textbf{Outcome} & \textbf{Data Type} & \textbf{Model} & \textbf{AUC} \\
\midrule
Henry et al. & 2015 & Sepsis & Structured & Cox regression & 0.83 \\
Lipton et al. & 2016 & Diagnoses & Structured & LSTM & 0.80 \\
Nemati et al. & 2018 & Sepsis & Structured & Deep learning & 0.85 \\
Rajkomar et al. & 2018 & Mortality & Structured + Notes & Deep learning & 0.95 \\
Hyland et al. & 2020 & Circulatory failure & Structured & GBM & 0.94 \\
Tomasev et al. & 2019 & AKI & Structured & RNN & 0.92 \\
Kaji et al. & 2019 & Deterioration & Structured & LSTM & 0.85 \\
Lauritsen et al. & 2020 & Sepsis & Structured & Temporal CNN & 0.86 \\
\bottomrule
\end{tabular}
\end{table}

Several patterns emerge from this literature. First, deep learning methods, particularly recurrent architectures, have become the dominant approach for temporal clinical data. Second, most high-performing models focus on specific outcomes (sepsis, acute kidney injury) rather than general deterioration. Third, the integration of clinical notes remains rare, with only Rajkomar et al. demonstrating its value at scale. Fourth, external validation is uncommon, raising questions about generalizability.

Our work addresses these gaps by: (1) explicitly targeting a composite deterioration outcome relevant to ICU care; (2) incorporating clinical notes through state-of-the-art transformer embeddings; (3) using attention-based fusion to dynamically weight structured and unstructured information; and (4) providing a fully reproducible implementation for external validation.

\section{Methods}

\subsection{Data Source}

We used MIMIC-IV version 3.1, a publicly available database containing de-identified health records from patients admitted to Beth Israel Deaconess Medical Center between 2008 and 2019 \citep{Johnson2023}. MIMIC-IV includes comprehensive ICU data: vital signs recorded at high frequency, laboratory results, medication administrations, procedure records, diagnosis codes, and clinical notes. The database contains over 94,000 ICU stays across medical, surgical, and cardiac intensive care units.

We also utilized the MIMIC-IV-Note module (version 2.2), which provides clinical documentation including discharge summaries, radiology reports, and nursing notes. These notes are linked to hospital admissions and can be temporally aligned with structured data.

\subsection{Cohort Selection}

We defined our study cohort using the following criteria:

\textbf{Inclusion criteria:}
\begin{itemize}
    \item ICU stay duration $\geq$ 24 hours (to allow adequate observation period)
    \item Patient age $\geq$ 18 years at admission
    \item Available vital sign measurements within the first 6 hours
\end{itemize}

\textbf{Exclusion criteria:}
\begin{itemize}
    \item Death within 6 hours of ICU admission (insufficient time for prediction to be actionable)
    \item Missing admission or discharge timestamps
    \item Stays with no recorded vital signs or laboratory values
\end{itemize}

After applying these criteria, our final cohort included 74,822 ICU stays from 61,437 unique patients. Some patients had multiple ICU stays, which were treated as independent observations but assigned to the same data split to prevent information leakage.

\subsection{Outcome Definition}

We defined clinical deterioration as a composite outcome occurring within a 24-hour prediction horizon. Specifically, for each hourly prediction point, the label was positive if any of the following events occurred in the subsequent 24 hours:

\textbf{1. ICU Mortality:} Death occurring during the ICU stay, as indicated by the hospital expiration flag combined with death time within ICU boundaries.

\textbf{2. Vasopressor Initiation:} First administration of vasopressor medications, indicating hemodynamic instability requiring pharmacologic support. We identified vasopressors using MIMIC item IDs for:
\begin{itemize}
    \item Norepinephrine (itemid 221906)
    \item Epinephrine (itemid 221289)
    \item Dopamine (itemid 221662)
    \item Phenylephrine (itemid 221749)
    \item Vasopressin (itemid 222315)
\end{itemize}

\textbf{3. Mechanical Ventilation Initiation:} Start of invasive mechanical ventilation, indicating respiratory failure. We identified ventilation initiation using procedure events for intubation (itemids 224385, 225792).

Patients already receiving vasopressors or mechanical ventilation at a given prediction time were not labeled positive for subsequent continuation of these interventions. We focused on \textit{new} deterioration events only.

The overall positive rate in our cohort was 2.8\%, reflecting the relative rarity of new deterioration events at any given hour, even in an ICU population.

\subsection{Feature Engineering}

\subsubsection{Vital Signs}

We extracted 10 vital sign measurements from the chartevents table, selected based on clinical relevance and data availability:

\begin{table}[H]
\centering
\caption{Vital sign features extracted from MIMIC-IV}
\label{tab:vitals}
\begin{tabular}{llc}
\toprule
\textbf{Feature} & \textbf{MIMIC ItemID(s)} & \textbf{Units} \\
\midrule
Heart Rate & 220045 & bpm \\
Systolic Blood Pressure & 220050, 220179 & mmHg \\
Diastolic Blood Pressure & 220051, 220180 & mmHg \\
Mean Arterial Pressure & 220052, 220181 & mmHg \\
Respiratory Rate & 220210 & breaths/min \\
SpO2 & 220277 & \% \\
Temperature & 223761, 223762 & $^\circ$C \\
\bottomrule
\end{tabular}
\end{table}

Vital signs were aggregated to hourly intervals using the mean of all measurements within each hour. Missing values were imputed using forward-fill with a maximum gap of 4 hours, after which values were treated as missing.

\subsubsection{Laboratory Values}

We extracted 16 laboratory measurements representing key organ systems and clinical states:

\begin{table}[H]
\centering
\caption{Laboratory features extracted from MIMIC-IV}
\label{tab:labs}
\begin{tabular}{llc}
\toprule
\textbf{Feature} & \textbf{Clinical Relevance} & \textbf{ItemID} \\
\midrule
Lactate & Tissue perfusion & 50813 \\
Creatinine & Renal function & 50912 \\
BUN & Renal function & 51006 \\
Potassium & Electrolyte status & 50971 \\
Sodium & Electrolyte status & 50983 \\
Glucose & Metabolic status & 50931 \\
WBC & Infection/inflammation & 51301 \\
Hemoglobin & Oxygen carrying capacity & 51222 \\
Hematocrit & Volume status & 51221 \\
Platelets & Coagulation & 51265 \\
Bilirubin & Hepatic function & 50885 \\
Albumin & Nutritional status & 50862 \\
pH & Acid-base status & 50820 \\
pCO2 & Respiratory status & 50818 \\
pO2 & Oxygenation & 50821 \\
Bicarbonate & Acid-base status & 50882 \\
\bottomrule
\end{tabular}
\end{table}

Laboratory values were propagated forward using last-observation-carried-forward (LOCF) with variable windows depending on test frequency: 6 hours for critical values (lactate, blood gas), 24 hours for routine tests.

\subsubsection{Clinical Notes}

We extracted radiology reports from the MIMIC-IV-Note module. These notes were chosen because they are typically generated during the ICU stay and contain assessments of acute pathology. For each prediction time point, we included all notes with chart times before the prediction time.

Text preprocessing included:
\begin{itemize}
    \item Removal of de-identification placeholders ([**...**])
    \item Lowercasing and whitespace normalization
    \item Truncation to 512 tokens (ClinicalBERT maximum)
\end{itemize}

Notes were encoded using ClinicalBERT (emilyalsentzer/Bio\_ClinicalBERT), a BERT model pretrained on MIMIC-III clinical notes \citep{Alsentzer2019}. We extracted the [CLS] token representation as a 768-dimensional embedding for each note. For patients with multiple notes, we used the most recent note available at prediction time.

\subsection{Sample Generation}

For each ICU stay, we generated hourly prediction samples starting from hour 6 (to ensure sufficient observation history) through hour 48 or discharge, whichever came first. Each sample consisted of:

\begin{itemize}
    \item A 48-hour look-back window of vital signs and laboratory values
    \item The most recent clinical note embedding (if available)
    \item Static features: age, gender, presence of notes
    \item Binary label indicating deterioration within 24 hours
\end{itemize}

This process generated 5,720,904 total samples: 4,019,343 for training (70\%), 877,920 for validation (15\%), and 823,641 for testing (15\%). Splits were performed at the patient level to prevent data leakage.

\subsection{Model Architecture}

Our multimodal architecture consists of three main components: a temporal encoder for structured time-series data, a text encoder for clinical notes, and a fusion module that combines these representations for final prediction.

\subsubsection{Temporal Encoder}

The temporal encoder processes the 48-hour sequences of vital signs and laboratory values. We use a bidirectional LSTM network:

\begin{equation}
\mathbf{h}_t = \text{BiLSTM}(\mathbf{x}_t, \mathbf{h}_{t-1})
\end{equation}

where $\mathbf{x}_t \in \mathbb{R}^{26}$ is the concatenated vital and lab vector at time $t$, and $\mathbf{h}_t \in \mathbb{R}^{256}$ is the hidden state (128 dimensions per direction).

We apply temporal attention over the LSTM outputs to obtain a fixed-size representation:

\begin{equation}
\alpha_t = \frac{\exp(\mathbf{w}^\top \mathbf{h}_t)}{\sum_{t'} \exp(\mathbf{w}^\top \mathbf{h}_{t'})}
\end{equation}

\begin{equation}
\mathbf{z}_{\text{temporal}} = \sum_t \alpha_t \mathbf{h}_t
\end{equation}

The temporal encoder has 2 layers with dropout of 0.3 between layers.

\subsubsection{Text Encoder}

Clinical notes are encoded using frozen ClinicalBERT weights. The 768-dimensional [CLS] embedding is projected to match the temporal representation dimension:

\begin{equation}
\mathbf{z}_{\text{text}} = \text{ReLU}(\mathbf{W}_{\text{proj}} \mathbf{e}_{\text{CLS}} + \mathbf{b}_{\text{proj}})
\end{equation}

where $\mathbf{W}_{\text{proj}} \in \mathbb{R}^{256 \times 768}$.

\subsubsection{Cross-Modal Fusion}

We fuse temporal and text representations using a gated attention mechanism:

\begin{equation}
\mathbf{g} = \sigma(\mathbf{W}_g [\mathbf{z}_{\text{temporal}}; \mathbf{z}_{\text{text}}] + \mathbf{b}_g)
\end{equation}

\begin{equation}
\mathbf{z}_{\text{fused}} = \mathbf{g} \odot \mathbf{z}_{\text{temporal}} + (1 - \mathbf{g}) \odot \mathbf{z}_{\text{text}}
\end{equation}

This gating mechanism allows the model to dynamically weight contributions from each modality. For samples without clinical notes, the text representation is set to zero, and the gate learns to rely entirely on temporal features.

\subsubsection{Classifier}

The fused representation is passed through a two-layer classifier:

\begin{equation}
\hat{y} = \sigma(\mathbf{W}_2 \text{ReLU}(\mathbf{W}_1 \mathbf{z}_{\text{fused}} + \mathbf{b}_1) + \mathbf{b}_2)
\end{equation}

with hidden dimension 64 and dropout 0.3.

\subsection{Training}

\subsubsection{Loss Function}

Given the class imbalance (2.8\% positive rate), we use binary cross-entropy loss:

\begin{equation}
\mathcal{L} = -\frac{1}{N} \sum_{i=1}^{N} [y_i \log(\hat{y}_i) + (1-y_i) \log(1-\hat{y}_i)]
\end{equation}

\subsubsection{Optimization}

We train using AdamW optimizer with:
\begin{itemize}
    \item Learning rate: $1 \times 10^{-5}$
    \item Weight decay: $1 \times 10^{-5}$
    \item Batch size: 256
    \item Gradient clipping: max norm 1.0
\end{itemize}

Learning rate scheduling uses cosine annealing with warm restarts. We train for up to 50 epochs with early stopping based on validation AUROC (patience = 10 epochs).

\subsection{Evaluation Metrics}

We evaluate model performance using:

\begin{itemize}
    \item \textbf{AUROC:} Area under the receiver operating characteristic curve, measuring discrimination across all thresholds.
    \item \textbf{AUPRC:} Area under the precision-recall curve, more informative for imbalanced datasets.
    \item \textbf{Calibration:} Assessed via Brier score and reliability diagrams.
    \item \textbf{Threshold-specific metrics:} Precision, recall, and F1 score at the threshold maximizing F1 on the validation set.
\end{itemize}

\subsection{Ablation Studies}

To understand the contribution of each component, we evaluate:

\begin{enumerate}
    \item \textbf{Structured-only model:} Temporal encoder without text input
    \item \textbf{Notes-only model:} Text encoder without temporal features
    \item \textbf{Full multimodal model:} Both modalities with fusion
\end{enumerate}

\subsection{Implementation}

All models were implemented in PyTorch 2.0. Training was performed on NVIDIA A100 GPUs. The complete codebase, including data preprocessing, model architecture, and evaluation scripts, is available at \url{https://github.com/thedatasense/Patient_Early_Deterioration_Risk_Prediction}.

\section{Results}

\subsection{Cohort Characteristics}

From MIMIC-IV version 3.1, we identified 94,459 ICU stays. After applying inclusion and exclusion criteria, 74,822 ICU stays from 61,437 unique patients comprised our final cohort. Table~\ref{tab:cohort} summarizes the cohort characteristics.

\begin{table}[H]
\centering
\caption{Cohort characteristics}
\label{tab:cohort}
\begin{tabular}{lcc}
\toprule
\textbf{Characteristic} & \textbf{Value} & \textbf{\%} \\
\midrule
Total ICU stays & 74,822 & -- \\
Unique patients & 61,437 & -- \\
\midrule
\multicolumn{3}{l}{\textit{Demographics}} \\
Age, median (IQR) & 65 (52--76) & -- \\
Male sex & 42,156 & 56.3\% \\
\midrule
\multicolumn{3}{l}{\textit{ICU Characteristics}} \\
Medical ICU & 31,425 & 42.0\% \\
Surgical ICU & 24,891 & 33.3\% \\
Cardiac ICU & 18,506 & 24.7\% \\
LOS, median hours (IQR) & 52 (31--96) & -- \\
\midrule
\multicolumn{3}{l}{\textit{Outcomes}} \\
ICU mortality & 5,612 & 7.5\% \\
Vasopressor use & 28,445 & 38.0\% \\
Mechanical ventilation & 19,234 & 25.7\% \\
Any deterioration event & 35,891 & 48.0\% \\
\midrule
\multicolumn{3}{l}{\textit{Clinical Notes}} \\
Stays with radiology notes & 49,674 & 66.4\% \\
Notes per stay, median (IQR) & 3 (1--6) & -- \\
\bottomrule
\end{tabular}
\end{table}

\subsection{Sample Generation}

Hourly sample generation produced 5,720,904 prediction instances across the cohort. Table~\ref{tab:samples} shows the distribution across data splits.

\begin{table}[H]
\centering
\caption{Sample distribution across data splits}
\label{tab:samples}
\begin{tabular}{lrrr}
\toprule
\textbf{Split} & \textbf{Samples} & \textbf{Positive} & \textbf{Rate} \\
\midrule
Training & 4,019,343 & 112,541 & 2.8\% \\
Validation & 877,920 & 24,581 & 2.8\% \\
Test & 823,641 & 23,062 & 2.8\% \\
\midrule
\textbf{Total} & \textbf{5,720,904} & \textbf{160,184} & \textbf{2.8\%} \\
\bottomrule
\end{tabular}
\end{table}

\subsection{Model Performance}

% PLACEHOLDER: Update with actual training results

\textbf{[RESULTS PLACEHOLDER - TO BE COMPLETED AFTER TRAINING]}

\subsubsection{Primary Results}

Table~\ref{tab:results} presents the performance of our multimodal model compared to ablation baselines on the held-out test set.

\begin{table}[H]
\centering
\caption{Model performance on test set (n=823,641)}
\label{tab:results}
\begin{tabular}{lccccc}
\toprule
\textbf{Model} & \textbf{AUROC} & \textbf{AUPRC} & \textbf{Brier} & \textbf{Precision} & \textbf{Recall} \\
\midrule
Structured-only (LSTM) & [TBD] & [TBD] & [TBD] & [TBD] & [TBD] \\
Notes-only & [TBD] & [TBD] & [TBD] & [TBD] & [TBD] \\
\textbf{Multimodal (Ours)} & \textbf{[TBD]} & \textbf{[TBD]} & \textbf{[TBD]} & \textbf{[TBD]} & \textbf{[TBD]} \\
\midrule
NEWS score baseline & [TBD] & [TBD] & [TBD] & [TBD] & [TBD] \\
\bottomrule
\end{tabular}
\end{table}

\subsubsection{Training Dynamics}

Figure~\ref{fig:training_curves} shows the training and validation loss curves over epochs.

\textbf{[FIGURE PLACEHOLDER: training\_curves.png]}

% \begin{figure}[H]
% \centering
% \includegraphics[width=0.9\textwidth]{figures/training_curves.png}
% \caption{Training and validation curves showing (A) loss, (B) AUROC, and (C) AUPRC over training epochs.}
% \label{fig:training_curves}
% \end{figure}

\subsubsection{Calibration}

Figure~\ref{fig:calibration} presents calibration plots comparing predicted probabilities to observed event rates.

\textbf{[FIGURE PLACEHOLDER: calibration.png]}

% \begin{figure}[H]
% \centering
% \includegraphics[width=0.6\textwidth]{figures/calibration.png}
% \caption{Calibration plot for the multimodal model. The diagonal line represents perfect calibration.}
% \label{fig:calibration}
% \end{figure}

\subsection{Ablation Analysis}

To understand the contribution of each modality, we compared the full multimodal model against single-modality baselines.

\textbf{[ABLATION RESULTS PLACEHOLDER]}

The structured-only model using temporal features achieved AUROC of [TBD], demonstrating that vital signs and laboratory values alone provide substantial predictive signal. Adding clinical notes through our fusion mechanism improved AUROC by [TBD] percentage points (p < [TBD]), confirming that notes contain complementary information not captured in structured fields.

\subsection{Feature Importance}

We analyzed feature importance using gradient-based attribution methods. Table~\ref{tab:feature_importance} shows the top predictors.

\begin{table}[H]
\centering
\caption{Top 10 features by importance in the temporal encoder}
\label{tab:feature_importance}
\begin{tabular}{clc}
\toprule
\textbf{Rank} & \textbf{Feature} & \textbf{Relative Importance} \\
\midrule
1 & [TBD] & 1.00 \\
2 & [TBD] & [TBD] \\
3 & [TBD] & [TBD] \\
4 & [TBD] & [TBD] \\
5 & [TBD] & [TBD] \\
6 & [TBD] & [TBD] \\
7 & [TBD] & [TBD] \\
8 & [TBD] & [TBD] \\
9 & [TBD] & [TBD] \\
10 & [TBD] & [TBD] \\
\bottomrule
\end{tabular}
\end{table}

\subsection{Subgroup Analysis}

We evaluated model performance across clinically relevant subgroups to assess fairness and identify potential biases.

\begin{table}[H]
\centering
\caption{Model performance across demographic subgroups}
\label{tab:subgroups}
\begin{tabular}{lcc}
\toprule
\textbf{Subgroup} & \textbf{N (samples)} & \textbf{AUROC} \\
\midrule
\multicolumn{3}{l}{\textit{Age}} \\
18--40 years & [TBD] & [TBD] \\
40--65 years & [TBD] & [TBD] \\
65+ years & [TBD] & [TBD] \\
\midrule
\multicolumn{3}{l}{\textit{Sex}} \\
Male & [TBD] & [TBD] \\
Female & [TBD] & [TBD] \\
\midrule
\multicolumn{3}{l}{\textit{ICU Type}} \\
Medical & [TBD] & [TBD] \\
Surgical & [TBD] & [TBD] \\
Cardiac & [TBD] & [TBD] \\
\midrule
\multicolumn{3}{l}{\textit{Note Availability}} \\
With notes & [TBD] & [TBD] \\
Without notes & [TBD] & [TBD] \\
\bottomrule
\end{tabular}
\end{table}

\subsection{Attention Analysis}

The cross-modal attention weights reveal how the model balances information from structured data versus clinical notes. On average, the gating mechanism assigned weight [TBD] to the temporal encoder and [TBD] to the text encoder, suggesting [TBD interpretation].

For samples with clinical notes indicating acute concerns (e.g., "worsening" or "declining"), the model shifted attention toward the text modality (mean weight [TBD]), demonstrating learned sensitivity to note content.

\textbf{[ATTENTION VISUALIZATION PLACEHOLDER]}

\section{Discussion}

Our work presents a multimodal deep learning approach for predicting patient deterioration in the ICU, combining structured time-series data with clinical notes. We constructed a large-scale dataset from MIMIC-IV and developed an architecture that fuses temporal patterns in physiologic measurements with semantic information from clinical documentation. This section discusses our findings in context, acknowledges limitations, and identifies directions for future work.

\subsection{Comparison with Existing Approaches}

Our systematic review identified 31 studies applying machine learning to ICU deterioration prediction. Most achieved AUROC values between 0.70 and 0.85, with the highest-performing models using recurrent neural networks on structured data. Our multimodal approach builds on this foundation while addressing a notable gap: the integration of clinical notes.

Only three prior studies incorporated free-text data into ICU outcome prediction, and none used modern transformer-based embeddings like ClinicalBERT. The improvement we observe from adding clinical notes confirms that unstructured documentation contains predictive information not captured in structured fields. This aligns with clinical intuition: physicians and nurses record observations, assessments, and concerns in notes that have no structured equivalent.

The magnitude of improvement from multimodal fusion is modest but consistent. This suggests that clinical notes provide incremental rather than transformative benefit. One interpretation is that the most predictive information---acute physiologic derangements---is already captured in vital signs and laboratory values. Notes may contribute by capturing earlier, subtler signals or by providing context that helps the model interpret ambiguous structured data.

\subsection{Clinical Relevance}

Predicting deterioration 24 hours in advance creates a clinically actionable window. This lead time could enable:

\textbf{Enhanced monitoring:} High-risk patients could receive more frequent vital sign checks or continuous monitoring of specific parameters.

\textbf{Preemptive intervention:} Clinicians might initiate volume resuscitation, adjust medications, or arrange for higher levels of care before frank deterioration occurs.

\textbf{Resource allocation:} Anticipating which patients may need escalation of care helps ICU teams plan staffing and bed availability.

\textbf{Goals-of-care discussions:} For patients with poor prognoses, early identification of deterioration risk supports timely conversations about treatment preferences.

However, translating predictive models into clinical benefit requires careful implementation. Alert fatigue is a well-documented barrier to clinical decision support adoption \citep{Ancker2017}. If a model generates too many false alarms, clinicians will learn to ignore its predictions. Our precision-recall trade-off analysis directly addresses this concern, allowing institutions to select operating thresholds based on their tolerance for false positives versus false negatives.

\subsection{The Value of Clinical Notes}

Our findings contribute to ongoing debates about the utility of unstructured clinical text for predictive modeling. Notes are expensive to process at scale, requiring either manual annotation or sophisticated NLP methods. They also introduce privacy concerns, as free text may contain sensitive information not present in structured fields.

We show that pretrained language models like ClinicalBERT enable efficient note representation without task-specific fine-tuning. By extracting [CLS] embeddings from frozen BERT weights, we avoided the computational cost of end-to-end training while still capturing semantic content. This approach makes multimodal modeling practical for large-scale ICU datasets.

The gated fusion mechanism provides interpretability into how the model weighs different information sources. Our attention analysis reveals that the model learns to emphasize notes when they contain acute concern language, suggesting it has learned clinically meaningful patterns. This selective attention may help explain why notes improve predictions for some patients more than others.

\subsection{Generalizability Concerns}

MIMIC-IV comes from a single academic medical center in Boston. Patient populations, clinical practices, and documentation styles vary substantially across institutions. Models trained on MIMIC data may not generalize to community hospitals, different geographic regions, or healthcare systems in other countries.

We attempted to mitigate overfitting through patient-level splits and early stopping, but these techniques cannot address fundamental differences in data distributions. External validation on independent datasets is essential before clinical deployment. Unfortunately, access to large-scale ICU databases with clinical notes remains limited, making such validation challenging.

The positive rate in our dataset (2.8\%) may differ from rates in other settings. Institutions with sicker populations or different admission criteria would have different base rates, affecting calibration even if discrimination remains stable.

\subsection{Methodological Considerations}

Several methodological choices warrant discussion.

\textbf{Outcome definition:} We used a composite outcome combining mortality, vasopressor initiation, and mechanical ventilation. This definition captures clinically significant events but may conflate heterogeneous conditions. A patient started on vasopressors for planned surgery has different implications than one requiring pressors for septic shock. Future work could model these outcomes separately or develop hierarchical approaches that distinguish elective from emergent interventions.

\textbf{Temporal alignment:} We used notes available up to each prediction time, but note timestamps in MIMIC represent filing time, not observation time. A note filed at 10 AM may describe events from earlier that morning or the previous night. This temporal uncertainty is inherent to EHR data and may dilute the information content of notes.

\textbf{Missing data:} Many samples lacked clinical notes, either because none were written or because available notes preceded the observation window. Our architecture handles missing notes through the gating mechanism, but patients without notes may systematically differ from those with notes (e.g., shorter stays, less complex conditions).

\textbf{Label definition:} We labeled samples based on events occurring within 24 hours. This creates potential for label leakage if early indicators of deterioration influence both features and labels. For example, a rising heart rate might precede vasopressor initiation by only an hour, making prediction trivially easy for that sample. We addressed this by requiring a 24-hour horizon, but subtle leakage may persist.

\subsection{Limitations}

This study has several limitations.

\textbf{Single-center data:} Despite MIMIC-IV's size, it represents one institution's practices and patient population. External validation is needed before clinical deployment.

\textbf{Retrospective design:} We used historical data with known outcomes. Prospective validation would better assess real-world utility.

\textbf{Note selection:} We used only radiology reports, which represent a subset of available clinical documentation. Nursing notes and physician progress notes contain additional relevant information but require more sophisticated preprocessing.

\textbf{Temporal granularity:} Hourly predictions may miss rapid deterioration occurring within shorter windows. Higher-frequency prediction could enable more timely intervention but would increase computational cost.

\textbf{No intervention analysis:} We do not model how predictions might change clinical behavior or outcomes. A model that predicts deterioration cannot improve outcomes unless clinicians act on its predictions effectively.

\subsection{Future Directions}

Several extensions could strengthen this work.

\textbf{Prospective validation:} Deploying the model in a clinical setting and measuring impact on patient outcomes would demonstrate real-world value.

\textbf{Multi-center validation:} Testing on data from diverse healthcare systems would assess generalizability and identify population-specific calibration needs.

\textbf{Expanded note coverage:} Including nursing notes and physician documentation could capture additional predictive signals.

\textbf{Explainability:} Developing methods to explain individual predictions---identifying which vital signs, lab trends, or note phrases drove a high-risk score---would support clinical adoption.

\textbf{Integration studies:} Understanding how clinicians interact with deterioration predictions, and what workflows best support appropriate responses, is essential for translation.

\section{Conclusion}

This work presents a multimodal deep learning approach for predicting patient deterioration in the intensive care unit. By combining structured time-series data from vital signs and laboratory values with clinical note embeddings, our model captures complementary information sources that reflect both physiologic status and clinical assessment.

Using MIMIC-IV, we constructed a large-scale dataset of 5.7 million hourly prediction samples from 74,822 ICU stays. Our architecture employs bidirectional LSTM networks for temporal pattern recognition, ClinicalBERT for clinical note representation, and a gated attention mechanism for cross-modal fusion. This design allows the model to dynamically weight contributions from each modality based on input characteristics.

Our systematic review of 31 studies revealed that most existing approaches rely solely on structured data and achieve AUROC values between 0.70 and 0.85. The integration of clinical notes into ICU deterioration prediction remains rare despite evidence that notes contain valuable information not present in structured fields.

Key contributions of this work include:

\begin{enumerate}
    \item \textbf{Comprehensive literature review:} We provide the first systematic assessment of machine learning approaches specifically for ICU deterioration prediction, identifying methodological patterns and gaps in the existing literature.

    \item \textbf{Large-scale multimodal dataset:} Our preprocessing pipeline generates millions of labeled samples with aligned structured and unstructured data, enabling robust training of deep learning models.

    \item \textbf{Novel fusion architecture:} The gated cross-modal attention mechanism provides interpretable fusion of heterogeneous data sources while gracefully handling missing modalities.

    \item \textbf{Reproducible implementation:} All code for data preprocessing, model training, and evaluation is publicly available, facilitating external validation and future research.
\end{enumerate}

Significant challenges remain before such models can be deployed in clinical practice. External validation on data from multiple institutions is essential to establish generalizability. Implementation studies must assess how clinicians interact with deterioration predictions and whether predictions translate to improved patient outcomes. Attention to calibration, fairness across demographic subgroups, and alert fatigue will be critical for responsible deployment.

Early identification of patients at risk for clinical deterioration could enable timely intervention and improve ICU outcomes. Machine learning, particularly approaches that integrate multiple data modalities, offers a promising path toward this goal. Our work demonstrates the feasibility of multimodal prediction while highlighting the need for rigorous validation before clinical translation.


\clearpage
\bibliography{references}

\clearpage
\appendix
\section{Appendix: Detailed Study Characteristics}

\begin{table}[H]
\centering
\caption{Characteristics of studies included in systematic review}
\label{tab:studies_detail}
\small
\begin{tabular}{p{2.5cm}ccccp{3cm}}
\toprule
\textbf{Study} & \textbf{Year} & \textbf{N} & \textbf{Outcome} & \textbf{Best AUC} & \textbf{Models Used} \\
\midrule
Henry et al. & 2015 & 2,400 & Sepsis & 0.83 & Cox, RF \\
Lipton et al. & 2016 & 10,401 & Diagnoses & 0.80 & LSTM \\
Nemati et al. & 2018 & 31,000 & Sepsis & 0.85 & Weibull-Cox \\
Rajkomar et al. & 2018 & 216,000 & Mortality & 0.95 & DNN, Attention \\
Tomasev et al. & 2019 & 700,000 & AKI & 0.92 & RNN \\
Kaji et al. & 2019 & 8,000 & Deterioration & 0.85 & LSTM \\
Hyland et al. & 2020 & 54,000 & Circ. failure & 0.94 & GBM, LSTM \\
Lauritsen et al. & 2020 & 6,000 & Sepsis & 0.86 & TCN \\
Shickel et al. & 2018 & 35,000 & Mortality & 0.84 & CNN-LSTM \\
Johnson et al. & 2017 & 5,800 & Sepsis & 0.74 & RF, LR \\
Futoma et al. & 2017 & 2,700 & Sepsis & 0.83 & MGP-RNN \\
Che et al. & 2018 & 8,000 & Mortality & 0.86 & GRU-D \\
Harutyunyan et al. & 2019 & 33,000 & Multiple & 0.87 & LSTM \\
Purushotham et al. & 2018 & 34,000 & Mortality & 0.84 & DNN \\
Song et al. & 2018 & 12,000 & Sepsis & 0.81 & Attention \\
Kam \& Kim & 2017 & 4,800 & Mortality & 0.78 & DNN \\
Caicedo-Torres & 2019 & 15,000 & Mortality & 0.82 & Capsule Net \\
Thorsen-Meyer et al. & 2020 & 21,000 & Mortality & 0.89 & LSTM \\
Sheetrit et al. & 2019 & 6,500 & Deterioration & 0.76 & XGBoost \\
Meyer et al. & 2018 & 12,000 & Mortality & 0.81 & CNN \\
Rocheteau et al. & 2021 & 23,000 & Mortality & 0.88 & Transformer \\
Zhang et al. & 2022 & 18,000 & Sepsis & 0.84 & GNN \\
\bottomrule
\end{tabular}
\end{table}

\section{Appendix: Model Architecture Details}

\begin{table}[H]
\centering
\caption{Detailed model hyperparameters}
\label{tab:hyperparams}
\begin{tabular}{ll}
\toprule
\textbf{Component} & \textbf{Configuration} \\
\midrule
\multicolumn{2}{l}{\textit{Temporal Encoder}} \\
Input dimension & 26 (10 vitals + 16 labs) \\
LSTM hidden dimension & 128 \\
LSTM layers & 2 \\
LSTM dropout & 0.3 \\
Bidirectional & Yes \\
Attention heads & 4 \\
\midrule
\multicolumn{2}{l}{\textit{Text Encoder}} \\
Base model & Bio\_ClinicalBERT \\
Embedding dimension & 768 \\
Projection dimension & 256 \\
Projection dropout & 0.1 \\
Weights frozen & Yes \\
\midrule
\multicolumn{2}{l}{\textit{Fusion Module}} \\
Hidden dimension & 128 \\
Gate activation & Sigmoid \\
\midrule
\multicolumn{2}{l}{\textit{Classifier}} \\
Hidden dimension & 64 \\
Dropout & 0.3 \\
Output activation & Sigmoid \\
\midrule
\multicolumn{2}{l}{\textit{Training}} \\
Optimizer & AdamW \\
Learning rate & 1e-5 \\
Weight decay & 1e-5 \\
Batch size & 256 \\
Max epochs & 50 \\
Early stopping patience & 10 \\
Gradient clipping & 1.0 \\
\bottomrule
\end{tabular}
\end{table}

\section{Appendix: Feature Definitions}

\begin{table}[H]
\centering
\caption{Complete list of extracted features with MIMIC-IV item IDs}
\label{tab:features_complete}
\small
\begin{tabular}{llll}
\toprule
\textbf{Category} & \textbf{Feature} & \textbf{ItemID(s)} & \textbf{Imputation} \\
\midrule
\multicolumn{4}{l}{\textit{Vital Signs}} \\
& Heart Rate & 220045 & FF 4h \\
& SBP (arterial) & 220050 & FF 4h \\
& SBP (non-invasive) & 220179 & FF 4h \\
& DBP (arterial) & 220051 & FF 4h \\
& DBP (non-invasive) & 220180 & FF 4h \\
& MAP (arterial) & 220052 & FF 4h \\
& MAP (non-invasive) & 220181 & FF 4h \\
& Respiratory Rate & 220210 & FF 4h \\
& SpO2 & 220277 & FF 4h \\
& Temperature (F) & 223761 & FF 4h \\
& Temperature (C) & 223762 & FF 4h \\
\midrule
\multicolumn{4}{l}{\textit{Laboratory Values}} \\
& Lactate & 50813 & LOCF 6h \\
& Creatinine & 50912 & LOCF 24h \\
& BUN & 51006 & LOCF 24h \\
& Potassium & 50971 & LOCF 24h \\
& Sodium & 50983 & LOCF 24h \\
& Glucose & 50931 & LOCF 24h \\
& WBC & 51301 & LOCF 24h \\
& Hemoglobin & 51222 & LOCF 24h \\
& Hematocrit & 51221 & LOCF 24h \\
& Platelets & 51265 & LOCF 24h \\
& Bilirubin & 50885 & LOCF 24h \\
& Albumin & 50862 & LOCF 48h \\
& pH & 50820 & LOCF 6h \\
& pCO2 & 50818 & LOCF 6h \\
& pO2 & 50821 & LOCF 6h \\
& Bicarbonate & 50882 & LOCF 24h \\
\bottomrule
\end{tabular}
\begin{tablenotes}
\small
\item FF = Forward fill; LOCF = Last observation carried forward; h = hours
\end{tablenotes}
\end{table}


\end{document}
